\documentclass[11pt,a4paper]{article}
\usepackage[utf8]{inputenc}
\usepackage[T1]{fontenc}
\usepackage[french]{babel}
\usepackage[margin=1.8cm,top=2cm,bottom=2cm]{geometry}
\usepackage{lmodern}
\usepackage{enumitem}
\usepackage{booktabs}
\usepackage{array}
\usepackage{longtable}
\usepackage{xcolor}
\usepackage{colortbl}
\usepackage{titlesec}
\usepackage{fancyhdr}
\usepackage{hyperref}
\usepackage{multicol}
\usepackage{graphicx}
\usepackage{tcolorbox}
\usepackage{tikz}
\usepackage{amssymb}

% Couleurs
\definecolor{afblue}{RGB}{0,48,135}
\definecolor{afred}{RGB}{200,16,46}
\definecolor{aflight}{RGB}{220,230,245}
\definecolor{memoyellow}{RGB}{255,250,220}
\definecolor{warnred}{RGB}{255,235,235}
\definecolor{okgreen}{RGB}{230,250,230}

% Style titres
\titleformat{\section}{\LARGE\bfseries\color{afblue}}{}{0em}{}[\titlerule]
\titleformat{\subsection}{\Large\bfseries\color{afblue!80}}{}{0em}{}
\titleformat{\subsubsection}{\large\bfseries\color{afblue!60}}{}{0em}{}

% Header/Footer
\pagestyle{fancy}
\fancyhf{}
\fancyhead[L]{\small\color{afblue}\textbf{PSY0 --- Cadets Air France}}
\fancyhead[R]{\small\color{afblue}Édition 2025/2026}
\fancyfoot[C]{\thepage}
\renewcommand{\headrulewidth}{1pt}
\renewcommand{\headrule}{\hbox to\headwidth{\color{afblue}\leaders\hrule height \headrulewidth\hfill}}

% Boîtes
\newtcolorbox{memobox}[1][]{colback=memoyellow,colframe=orange!70,title=\textbf{MÉMO},fonttitle=\bfseries,#1}
\newtcolorbox{warnbox}[1][]{colback=warnred,colframe=red!70,title=\textbf{ATTENTION},fonttitle=\bfseries,#1}
\newtcolorbox{tipbox}[1][]{colback=okgreen,colframe=green!50!black,title=\textbf{ASTUCE},fonttitle=\bfseries,#1}

\hypersetup{colorlinks=true,linkcolor=afblue,urlcolor=afblue}

\begin{document}

% PAGE DE TITRE
\begin{titlepage}
\begin{center}
\vspace*{2cm}
{\Huge\bfseries\color{afblue} GUIDE DE RÉVISION\\[0.3cm] ULTRA-COMPLET}\\[1cm]
{\LARGE\color{afred}\bfseries PSY0 --- CADETS AIR FRANCE}\\[0.8cm]
{\Large Culture Aéronautique --- Édition 2025/2026}\\[2cm]
\rule{\textwidth}{2pt}\\[0.5cm]
{\large Ce document compile l'intégralité des connaissances nécessaires pour réussir l'épreuve de culture aéronautique du PSY0.\\[0.3cm]
Basé sur les annales 2018, 2019, 2020, les fiches de préparation officielles et des recherches complémentaires actualisées.}\\[0.5cm]
\rule{\textwidth}{2pt}\\[2cm]
{\Large\textbf{15 Parties $\bullet$ 100 QCM $\bullet$ Fiches Mémo}}\\[3cm]
{\large\color{gray} Bonne chance pour ta sélection, futur(e) pilote !}
\end{center}
\end{titlepage}

\tableofcontents
\newpage

%=============================================================
\section{HISTOIRE DE L'AVIATION}
%=============================================================

\subsection{Les Pionniers (avant 1914)}

\begin{longtable}{p{1cm}p{3.5cm}p{4cm}p{6cm}}
\toprule
\textbf{Année} & \textbf{Qui} & \textbf{Exploit} & \textbf{Détail à retenir} \\
\midrule
\endfirsthead
\toprule \textbf{Année} & \textbf{Qui} & \textbf{Exploit} & \textbf{Détail} \\ \midrule
\endhead
1783 & Frères Montgolfier & Premier vol en ballon & Ballon à air chaud, Annonay (Ardèche) \\
1783 & Jacques Charles & Ballon à gaz (hydrogène) & Tuileries, Paris \\
1852 & Henri Giffard & Dirigeable motorisé & Machine à vapeur \\
\rowcolor{aflight}
1890 & \textbf{Clément Ader} & \textbf{Premier décollage motorisé (Éole)} & $\sim$50m, non contrôlé. <<~Père de l'aviation~>> en France. Invente le mot <<~avion~>> \\
1891 & Otto Lilienthal & Vols planés répétés & Allemand, $\sim$2000 vols planés \\
\rowcolor{aflight}
1903 & \textbf{Frères Wright} & \textbf{Premier vol motorisé contrôlé} & Kitty Hawk (USA), 12s, 36m, \textit{Wright Flyer} \\
1906 & Santos-Dumont & Vol du 14-bis à Bagatelle & Premier vol homologué en Europe \\
\rowcolor{aflight}
1909 & \textbf{Louis Blériot} & \textbf{Traversée de la Manche} & Blériot XI, 25 juillet, 37 min \\
1910 & Henri Fabre & Premier vol en hydravion & Étang de Berre (Marseille) \\
1910 & Élisa Deroche & Première femme brevetée pilote & Brevet n°36 (monde entier) \\
1911 & Raoul Badin & Invente l'anémomètre & Tube de Badin --- mesure vitesse air \\
\rowcolor{aflight}
1913 & \textbf{Roland Garros} & \textbf{Traversée de la Méditerranée} & Fréjus $\to$ Bizerte (Tunisie), 8h \\
1913 & Adolphe Pégoud & Premier looping en avion & Français \\
\bottomrule
\end{longtable}

\subsection{Grande Guerre et Entre-deux-guerres (1914--1939)}

\begin{longtable}{p{1cm}p{4cm}p{9.5cm}}
\toprule
\textbf{Année} & \textbf{Événement} & \textbf{Détails} \\
\midrule
1914-18 & WWI & Reconnaissance $\to$ chasse $\to$ bombardement. As~: Guynemer (54 vict.), Fonck (75 vict., top as allié), Richthofen <<~Baron Rouge~>> (80 vict.) \\
1919 & Création de KLM & Plus ancienne compagnie encore en activité \\
1921 & Adrienne Bolland & \textbf{Première femme à traverser les Andes} (Caudron G.3) \\
\rowcolor{aflight}
1927 & \textbf{Charles Lindbergh} & \textbf{Première traversée solo Atlantique Nord sans escale}. NY $\to$ Paris (Le Bourget), \textit{Spirit of St. Louis}, 33h30 \\
1927 & Nungesser et Coli & Tentative Paris $\to$ NY, disparus. Avion~: <<~L'Oiseau Blanc~>> \\
1930 & Jean Mermoz & Première traversée postale Atlantique Sud (Latécoère 28) \\
\rowcolor{aflight}
\textbf{1933} & \textbf{Création d'AIR FRANCE} & Fusion de \textbf{5 compagnies}~: Air Orient, Air Union, CIDNA, Lignes Farman (SGTA), Aéropostale \\
1936 & Disparition de Mermoz & Latécoère 300 <<~Croix du Sud~>>, Atlantique Sud \\
\bottomrule
\end{longtable}

\subsection{WWII et Âge du Jet (1939--1970)}

\begin{longtable}{p{1cm}p{4cm}p{9.5cm}}
\toprule
\textbf{Année} & \textbf{Événement} & \textbf{Détails} \\
\midrule
1939-45 & WWII & Bond techno~: radar, réacteur, pressurisation. Spitfire, P-51 Mustang, Messerschmitt Bf 109 \\
1939 & Premier vol à réaction & Heinkel He 178 (turboréacteur Von Ohain) \\
\rowcolor{aflight}
1944 & \textbf{Convention de Chicago} & Création de l'\textbf{OACI} \\
\rowcolor{aflight}
1947 & \textbf{Chuck Yeager} & \textbf{Premier vol supersonique} --- Bell X-1, Mach 1.06 \\
1949 & De Havilland Comet & Premier avion de ligne à réaction \\
1955 & Caravelle --- premier vol & Premier biréacteur commercial de série, moteurs arrière \\
\rowcolor{aflight}
1969 & \textbf{Concorde} --- premier vol & 2 mars 1969, Toulouse, franco-britannique \\
1969 & Boeing 747 --- premier vol & <<~Jumbo Jet~>>, premier gros porteur \\
\rowcolor{aflight}
1970 & \textbf{Création d'Airbus} & 18 décembre, consortium européen \\
\bottomrule
\end{longtable}

\subsection{Ère Moderne (1970--aujourd'hui)}

\begin{longtable}{p{1cm}p{4cm}p{9.5cm}}
\toprule
\textbf{Année} & \textbf{Événement} & \textbf{Détails} \\
\midrule
1976 & Concorde en service AF+BA & Paris-NY 3h30, Mach 2.02, plafond 18\,000m \\
1988 & A320 en service chez AF & Premier avion \textbf{fly-by-wire} civil Airbus \\
\rowcolor{aflight}
1996 & \textbf{Hub CDG} inauguré & Sous Christian Blanc \\
\rowcolor{aflight}
2000 & Alliance \textbf{SkyTeam} fondée & AF, Delta, AeroMéxico, Korean Air \\
2003 & Arrêt du Concorde & Dernier vol~: 26 novembre 2003 \\
\rowcolor{aflight}
2004 & Groupe \textbf{Air France-KLM} & Fusion \\
2005 & Privatisation d'AF & + Lancement Flying Blue \\
2006 & Création Transavia France & Filiale low-cost \\
\rowcolor{aflight}
2019 & \textbf{A350-900 chez AF} & Remplace A340/A330/B777-200 \\
\rowcolor{aflight}
2020 & \textbf{Retrait A380} flotte AF & Covid accélère la décision \\
\bottomrule
\end{longtable}

\begin{memobox}
\textbf{Séquence chronologique}~: 1890 Ader $\to$ 1903 Wright $\to$ 1909 Blériot $\to$ 1927 Lindbergh $\to$ \textbf{1933 Air France} $\to$ 1947 Mur du son $\to$ 1969 Concorde $\to$ 1996 Hub CDG $\to$ 2000 SkyTeam $\to$ 2004 AF-KLM $\to$ 2019 A350 AF
\end{memobox}

%=============================================================
\section{AIR FRANCE --- IDENTITÉ, GOUVERNANCE, STRATÉGIE}
%=============================================================

\subsection{Fiche d'Identité}

\begin{tabular}{p{4.5cm}p{10cm}}
\toprule
Date de création & \textbf{7 octobre 1933} \\
Fusion de & Air Orient, Air Union, CIDNA, Lignes Farman (SGTA), Aéropostale \\
Nationalisation / Privatisation & 1945 / 2005 \\
Siège social & Tremblay-en-France (Roissy-CDG) \\
Code OACI / IATA & AFR / AF \\
Emblème & \textbf{Hippocampe ailé} (<<~la crevette~>>), hérité d'Air Orient \\
Slogan actuel (2022) & <<~S'envoler en toute élégance~>> \\
Alliance & SkyTeam (membre fondateur) \\
Programme fidélité & Flying Blue \\
\bottomrule
\end{tabular}

\subsection{Slogans Historiques}
\begin{itemize}[nosep]
\item 1980~: <<~Gagner le c\oe ur du monde~>>
\item 1999~: <<~Faire du ciel le plus bel endroit de la terre~>>
\item 2014--2022~: <<~\textbf{France is in the Air}~>> (agence BETC)
\item Depuis 2022~: <<~\textbf{S'envoler en toute élégance}~>>
\end{itemize}

\subsection{Dirigeants Clés}

\begin{tabular}{p{4cm}p{2.5cm}p{8cm}}
\toprule
\textbf{Nom} & \textbf{Période} & \textbf{Faits marquants} \\
\midrule
Ernest Roume & 1933--1935 & Premier président d'AF \\
Christian Blanc & 1993--1997 & Création du Hub CDG (1996) \\
Jean-Cyril Spinetta & 1997--2009 & SkyTeam, fusion KLM, privatisation, Flying Blue \\
\rowcolor{aflight}
\textbf{Benjamin Smith} & \textbf{Depuis 2018} & \textbf{DG Groupe AF-KLM} \\
\rowcolor{aflight}
\textbf{Anne Rigail} & \textbf{Depuis 2018} & \textbf{DG Air France} (première femme au poste) \\
Anne-Marie Couderc & Actuelle & Présidente du CA (non exécutif) \\
\bottomrule
\end{tabular}

\subsection{Filiales et Chiffres Clés}

\begin{multicols}{2}
\textbf{Filiales~:}
\begin{itemize}[nosep]
\item Air France (long + moyen courrier)
\item KLM (Pays-Bas, hub Amsterdam)
\item Transavia France (low-cost, hub Orly)
\item Air France HOP (régional)
\end{itemize}
\columnbreak
\textbf{Chiffres~:}
\begin{itemize}[nosep]
\item Flotte AF~: $\sim$\textbf{228 avions}
\item Destinations~: $\sim$200 dans $\sim$90 pays
\item Passagers/an AF~: $\sim$50 millions
\item Employés AF~: $\sim$43\,000
\end{itemize}
\end{multicols}

%=============================================================
\section{FLOTTE DÉTAILLÉE}
%=============================================================

\subsection{Flotte Long-Courrier}

\begin{tabular}{p{3cm}p{1.5cm}p{3cm}p{1.5cm}p{1.5cm}p{1.5cm}p{1.5cm}}
\toprule
\textbf{Avion} & \textbf{Nb} & \textbf{Moteur} & \textbf{Pax} & \textbf{Auton.} & \textbf{Mach} & \textbf{MTOW} \\
\midrule
\rowcolor{aflight}
\textbf{A350-900} & $\sim$40 & 2$\times$ Trent XWB (RR) & 324 & 15\,000 km & 0.85 & 280t \\
\textbf{B777-300ER} & $\sim$43 & 2$\times$ GE90-115B & 468 & 13\,650 km & 0.84 & 351t \\
B777-200ER & $\sim$18 & 2$\times$ GE90-94B & 280 & 14\,300 km & 0.84 & 298t \\
\bottomrule
\end{tabular}

\begin{tipbox}
Le \textbf{GE90-115B} du B777-300ER est le \textbf{moteur le plus puissant au monde}~: 115\,000 lb de poussée.
\end{tipbox}

\subsection{Flotte Moyen-Courrier}

\begin{tabular}{p{3cm}p{1.5cm}p{3cm}p{1.5cm}p{1.5cm}p{1.5cm}p{1.5cm}}
\toprule
\textbf{Avion} & \textbf{Nb} & \textbf{Moteur} & \textbf{Pax} & \textbf{Auton.} & \textbf{Mach} & \textbf{MTOW} \\
\midrule
\rowcolor{aflight}
\textbf{A220-300} & $\sim$50 & 2$\times$ PW1500G (P\&W) & 148 & 6\,300 km & 0.82 & 70.9t \\
A320 & actif & 2$\times$ CFM56-5 & 150 & 6\,150 km & 0.78 & 78t \\
A321 & actif & 2$\times$ CFM56-5 & 185 & 5\,950 km & 0.78 & 93.5t \\
\bottomrule
\end{tabular}

\subsection{Transavia France}
B737-800 ($\sim$33, 2$\times$ CFM56-7B, 189 pax) + A320neo en livraison (2$\times$ LEAP-1A).

\subsection{Avions Historiques}

\begin{tabular}{p{3cm}p{2cm}p{9.5cm}}
\toprule
\textbf{Avion} & \textbf{Dates AF} & \textbf{Caractéristiques} \\
\midrule
\textbf{Concorde} & 1976--2003 & \textbf{Mach 2.02}, plafond 18\,300m, 100 pax, 4$\times$ Olympus 593 \\
\textbf{A380} & 2009--2020 & 516 pax AF, 4$\times$ GP7200, double pont intégral \\
Caravelle & 1959--1981 & Premier biréacteur commercial série, moteurs arrière \\
Boeing 747 & 1970--2016 & <<~Jumbo Jet~>>, 4 moteurs, $\sim$400 pax \\
A340 & 1993--2020 & \textbf{4$\times$ CFM56}, quadriréacteur \\
\bottomrule
\end{tabular}

\subsection{Moteurs par Avion (question très fréquente)}

\begin{memobox}
\textbf{CFM International} = joint-venture \textbf{GE (USA) + Safran (France)}. Produit le CFM56 ($>$35\,000 ex.) et le LEAP.\\
\textbf{GE90} $\to$ B777 \quad|\quad \textbf{Trent XWB} $\to$ A350 \quad|\quad \textbf{CFM56} $\to$ A320/A340/B737 \quad|\quad \textbf{PW1500G} $\to$ A220
\end{memobox}

%=============================================================
\section{PRINCIPES DE VOL \& AÉRODYNAMIQUE}
%=============================================================

\subsection{Les 4 Forces}
En vol stabilisé en palier~: Portance = Poids \textbf{ET} Poussée = Traînée.

\begin{itemize}[nosep]
\item \textbf{Portance}~: Différence de pression extrados/intrados. Dépend de $V^2$, surface alaire, $C_z$, $\rho$.
\item \textbf{Poids}~: $m \times g$ ($g = 9{,}81$ m/s$^2$).
\item \textbf{Poussée}~: Force motrice (réacteurs/hélices).
\item \textbf{Traînée}~: Résistance aéro = traînée de forme + induite + frottement.
\end{itemize}

\subsection{Les 3 Axes et Gouvernes}

\begin{tabular}{p{2.5cm}p{2.5cm}p{2.5cm}p{3cm}p{3.5cm}}
\toprule
\textbf{Axe} & \textbf{Mouvement} & \textbf{Commande} & \textbf{Gouverne} & \textbf{Mnémo} \\
\midrule
Longitudinal & \textbf{Roulis} & Manche latéral & \textbf{Ailerons} & Roulis = Roll = Ailerons \\
Latéral & \textbf{Tangage} & Manche avant/arr. & \textbf{Profondeur} & Tangage = Pitch = Profondeur \\
Vertical & \textbf{Lacet} & Palonnier (pédales) & \textbf{Direction} & Lacet = Yaw = Direction \\
\bottomrule
\end{tabular}

\subsection{Relation Fondamentale}
$$\boxed{\text{ASSIETTE} = \text{PENTE} + \text{INCIDENCE}}$$

\subsection{Le Décrochage (Stall)}

\begin{warnbox}
Le décrochage se produit \textbf{toujours à la même incidence critique} ($\sim$15--18°), \textbf{indépendamment} de la vitesse, du poids ou de l'altitude. Par contre, la \textbf{vitesse} de décrochage varie avec le poids, le facteur de charge et la configuration des volets.
\end{warnbox}

\subsection{Dispositifs Hypersustentateurs}

\begin{tabular}{p{3cm}p{2.5cm}p{3cm}p{5.5cm}}
\toprule
\textbf{Dispositif} & \textbf{Position} & \textbf{Rôle} & \textbf{Détails} \\
\midrule
\textbf{Volets} (flaps) & Bord de fuite & $\uparrow$ portance, $\uparrow$ traînée & Décollage $\sim$10--15°, atterrissage $\sim$30--40° \\
\textbf{Becs} (slats) & Bord d'attaque & $\uparrow$ incidence max & Retardent le décrochage \\
\textbf{Spoilers} & Extrados & $\downarrow$ portance, $\uparrow$ traînée & Descente + freinage \\
\textbf{Reverses} & Nacelles moteurs & Inversent la poussée & Uniquement après toucher des roues \\
\bottomrule
\end{tabular}

%=============================================================
\section{MOTEURS \& PROPULSION}
%=============================================================

\subsection{Fonctionnement du Turboréacteur (Turbofan)}

\begin{center}
\fbox{\textbf{Fan $\to$ Compresseur $\to$ Chambre de Combustion $\to$ Turbine $\to$ Tuyère}}
\end{center}

\begin{memobox}
\textbf{Mnémo}~: \textbf{C}ompresseur $\to$ \textbf{C}ombustion $\to$ \textbf{T}urbine $\to$ \textbf{T}uyère\\
\textbf{Taux de dilution}~: CFM56 $\sim$6 | LEAP $\sim$11 | Militaire $\sim$0.5. Plus élevé = plus silencieux et économique.
\end{memobox}

\subsection{Types de Moteurs}

\begin{tabular}{p{3.5cm}p{5cm}p{2.5cm}p{3cm}}
\toprule
\textbf{Type} & \textbf{Principe} & \textbf{Vitesse opt.} & \textbf{Exemples} \\
\midrule
Turbofan (double flux) & Fan + compresseur + combustion & M 0.7--0.9 & CFM56, GE90, LEAP \\
Turbopropulseur & Turbine $\to$ hélice via réducteur & $<$ 600 km/h & PW100 (ATR) \\
Turboréacteur pur & Tout en flux primaire & $>$ Mach 1 & Olympus 593 \\
Moteur à pistons & Combustion interne + hélice & $<$ 400 km/h & Lycoming, Continental \\
\bottomrule
\end{tabular}

%=============================================================
\section{INSTRUMENTS DE BORD}
%=============================================================

\subsection{Les Six Instruments de Base}

\begin{tabular}{p{3cm}p{3cm}p{4cm}p{4.5cm}}
\toprule
\textbf{Instrument} & \textbf{Mesure} & \textbf{Source} & \textbf{Mnémo} \\
\midrule
Anémomètre (Badin) & Vitesse (IAS) & Pitot $-$ Statique & -- \\
Horizon artificiel & Assiette & Gyroscope / AHRS & Bleu=ciel, Marron=terre \\
Altimètre & Altitude & Pression statique & QNH=Alt, QFE=Haut, 1013=FL \\
Variomètre & Vit. verticale & Variation pression stat. & ft/min, + = montée \\
Conservateur de cap & Cap & Gyroscope directionnel & -- \\
Bille-aiguille & Virage + symétrie & Gyro + pendule & Bille au centre \\
\bottomrule
\end{tabular}

\subsection{Pannes --- Effets (question très fréquente)}

\begin{warnbox}[title=\textbf{Si la prise STATIQUE est bouchée}]
\textbf{Altimètre}~: bloqué à la dernière altitude.
\textbf{Variomètre}~: bloqué à zéro.
\textbf{Badin}~: sous-estime en montée, surestime en descente.
\end{warnbox}

\begin{warnbox}[title=\textbf{Si le tube PITOT est bouché}]
\textbf{Badin}~: indications erronées (se comporte comme un altimètre). Altimètre et variomètre restent OK.
\end{warnbox}

\subsection{Systèmes de Sécurité}
\begin{itemize}[nosep]
\item \textbf{TCAS}~: anticollision (TA/RA), basé sur transpondeurs
\item \textbf{GPWS/EGPWS}~: alerte proximité sol (<<~TERRAIN PULL UP~>>)
\item \textbf{FMS}~: ordinateur navigation/performances
\item \textbf{ECAM} (Airbus) / \textbf{EICAS} (Boeing)~: affichage pannes
\item \textbf{CVR}~: boîte noire voix ($\geq$2h) | \textbf{FDR}~: boîte noire données ($\geq$25h)
\end{itemize}

%=============================================================
\section{VITESSES AÉRONAUTIQUES}
%=============================================================

\subsection{Vitesses de Décollage}

\begin{tabular}{p{2cm}p{4cm}p{8.5cm}}
\toprule
\textbf{Vitesse} & \textbf{Nom} & \textbf{Définition} \\
\midrule
\textbf{V1} & Vitesse de décision & Max pour interrompre le décollage. Au-delà $\to$ on continue \\
\textbf{VR} & Vitesse de rotation & On tire le manche (VR $\geq$ V1) \\
\textbf{V2} & Vitesse de sécurité & Min pour monter avec 1 moteur en panne \\
\bottomrule
\end{tabular}

\subsection{Types de Vitesses}

\begin{tabular}{p{2cm}p{5cm}p{7.5cm}}
\toprule
\textbf{Abrév.} & \textbf{Nom} & \textbf{Explication} \\
\midrule
\textbf{IAS} & Indicated Airspeed & Lue sur le Badin. Max 250 kt sous FL100 \\
\textbf{TAS} & True Airspeed & Vitesse vraie / air. Augmente avec l'altitude (à même IAS) \\
\textbf{GS} & Ground Speed & TAS $\pm$ vent. Détermine le temps de vol réel \\
\textbf{Mach} & Nombre de Mach & TAS / vitesse du son locale. Utilisé au-dessus du FL250 \\
\bottomrule
\end{tabular}

\subsection{Vitesses Typiques A320}

\begin{tabular}{p{5cm}p{9.5cm}}
\toprule
\textbf{Paramètre} & \textbf{Valeur} \\
\midrule
V1 / VR / V2 (typ.) & $\sim$145 / $\sim$150 / $\sim$155 kt \\
\rowcolor{aflight}
\textbf{Croisière} & \textbf{Mach 0.78} (\textbf{828 km/h}) \\
MMO / VMO & Mach 0.82 / 350 kt IAS \\
Altitude croisière & FL350 à FL390 (plafond FL410) \\
Consommation horaire & $\sim$2\,700 kg/h par moteur ($\sim$5\,400 L/h total) \\
\bottomrule
\end{tabular}

\newpage
%=============================================================
\section{NAVIGATION AÉRIENNE}
%=============================================================

\subsection{Systèmes de Navigation}

\begin{tabular}{p{1.8cm}p{2cm}p{1.5cm}p{3cm}p{6cm}}
\toprule
\textbf{Système} & \textbf{Bande} & \textbf{Portée} & \textbf{Infos} & \textbf{Particularités} \\
\midrule
\textbf{NDB} & MF & $\sim$200 NM & Direction (ADF) & Peu précis, sensible parasites/orages \\
\textbf{VOR} & VHF & $\sim$200 NM & Radiale (relèvement) & Précis, portée optique \\
\textbf{DME} & UHF & $\sim$200 NM & Distance oblique (NM) & Souvent couplé au VOR \\
\textbf{ILS} & VHF+UHF & $\sim$25 NM & Axe piste + pente 3° & Le plus précis pour atterrir \\
\textbf{GNSS} & Satellite & Mondiale & Position 3D & GPS/Galileo/GLONASS/BeiDou \\
\textbf{IRS/INS} & Embarqué & Autonome & Position, vitesse & Dérive $\sim$1 NM/h \\
\bottomrule
\end{tabular}

\subsection{Catégories ILS}

\begin{tabular}{p{2cm}p{5cm}p{7.5cm}}
\toprule
\textbf{Cat.} & \textbf{DH (Decision Height)} & \textbf{RVR (Runway Visual Range)} \\
\midrule
Cat I & $\geq$ 200 ft & $\geq$ 550 m \\
Cat II & $\geq$ 100 ft & $\geq$ 300 m \\
Cat IIIA & $\geq$ 50 ft & $\geq$ 200 m \\
Cat IIIB & $<$ 50 ft & $\geq$ 75 m \\
Cat IIIC & Aucun minimum & Aucun minimum \\
\bottomrule
\end{tabular}

\subsection{Espaces Aériens en France}

\begin{tabular}{p{1.5cm}p{2cm}p{1cm}p{1cm}p{3.5cm}p{2cm}p{2cm}}
\toprule
\textbf{Classe} & \textbf{Type} & \textbf{IFR} & \textbf{VFR} & \textbf{Séparation} & \textbf{Radio} & \textbf{Clair.} \\
\midrule
A & Contrôlé & Oui & Non & IFR/IFR & Oblig. & Oblig. \\
B & --- & \multicolumn{5}{c}{\textit{Non utilisé en France}} \\
C & Contrôlé & Oui & Oui & IFR/IFR + IFR/VFR & Oblig. & Oblig. \\
D & Contrôlé & Oui & Oui & IFR/IFR seulement & Oblig. & Oblig. \\
E & Contrôlé & Oui & Oui & IFR/IFR & IFR: oui & IFR seul \\
F & --- & \multicolumn{5}{c}{\textit{Non utilisé en France}} \\
G & Non contrôlé & Oui & Oui & Aucune & Non oblig. & Aucune \\
\bottomrule
\end{tabular}

\subsection{Niveaux de Vol --- Règle Semi-Circulaire}
\begin{itemize}[nosep]
\item Route 000°--179° (Est)~: FL \textbf{impairs} (310, 330, 350, 370...)
\item Route 180°--359° (Ouest)~: FL \textbf{pairs} (320, 340, 360, 380...)
\item Au-dessus FL290~: espacement RVSM de 1\,000 ft
\end{itemize}

\subsection{Navigation Transatlantique}
Pas de VOR/radar au-dessus de l'océan $\to$ \textbf{GPS + IRS}. Routes NAT tracks (révisées 2$\times$/jour selon jet stream).

%=============================================================
\section{MÉTÉOROLOGIE AÉRONAUTIQUE}
%=============================================================

\subsection{Atmosphère ISA}

\begin{tabular}{p{4cm}p{10.5cm}}
\toprule
Température au sol & \textbf{15°C} \\
Pression au sol & \textbf{1013.25 hPa} \\
Gradient thermique & \textbf{$-$2°C / 1\,000 ft} ($-$6.5°C/km) \\
Tropopause & $\sim$36\,000 ft (11 km), T = $-$56.5°C \\
\bottomrule
\end{tabular}

\subsection{Tableau Températures ISA}

\begin{tabular}{p{3cm}p{3cm}p{3cm}p{3cm}p{3cm}}
\toprule
0 ft & 5\,000 ft & 10\,000 ft & 20\,000 ft & 36\,000 ft \\
\midrule
\textbf{+15°C} & \textbf{+5°C} & \textbf{$-$5°C} & \textbf{$-$25°C} & \textbf{$-$56.5°C} \\
\bottomrule
\end{tabular}

\subsection{Calages Altimétriques}

\begin{tabular}{p{3cm}p{4cm}p{7.5cm}}
\toprule
\textbf{QNH} & Affiche l'\textbf{altitude} & En dessous de l'altitude de transition \\
\textbf{QFE} & Affiche la \textbf{hauteur} & Au-dessus du terrain (peu utilisé en France) \\
\textbf{1013.25 hPa} & Affiche le \textbf{niveau de vol} & Au-dessus de l'altitude de transition \\
\bottomrule
\end{tabular}

\begin{memobox}
1 hPa $\approx$ 28 ft. \textbf{<<~High to low, look out below!~>>} (haute $\to$ basse pression = on descend sans le savoir)
\end{memobox}

\subsection{Nuages}

\begin{tabular}{p{2cm}p{3cm}p{2cm}p{7cm}}
\toprule
\textbf{Étage} & \textbf{Type} & \textbf{Altitude} & \textbf{Caractéristiques} \\
\midrule
Haut & Cirrus, Cs, Cc & 6--12 km & Cristaux de glace \\
Moyen & Altostratus, Ac & 2--6 km & Voile gris \\
Bas & Stratus, Sc, Ns & 0--2 km & Brume, pluie continue \\
\rowcolor{warnred}
Dév. vert. & \textbf{Cumulonimbus (Cb)} & 0.5--15 km & \textbf{LE PLUS DANGEREUX}~: orages, grêle, cisaillement \\
\bottomrule
\end{tabular}

\subsection{Phénomènes Clés}

\begin{itemize}[nosep]
\item \textbf{Jet stream}~: 150--300 kt à la tropopause, d'ouest en est $\to$ NY$\to$Paris plus rapide
\item \textbf{Anticyclone (H)}~: haute pression, sens horaire (HN), beau temps
\item \textbf{Dépression (L)}~: basse pression, sens antihoraire (HN), mauvais temps
\item \textbf{Cisaillement de vent}~: changement brutal, très dangereux basse altitude
\item \textbf{Givrage}~: glace sur cellule/sondes (cause crash AF447)
\end{itemize}

\subsection{METAR --- Décodage}

Exemple~: \texttt{METAR LFPG 101330Z 25015G28KT 9999 FEW035 SCT080 18/09 Q1018}

\begin{tabular}{p{3cm}p{11.5cm}}
\toprule
LFPG & Paris CDG \\
101330Z & Le 10, 13h30 UTC \\
25015G28KT & Vent 250° à 15 kt, rafales 28 kt \\
9999 & Visibilité $>$ 10 km \\
FEW035 SCT080 & Quelques nuages 3\,500 ft, épars 8\,000 ft \\
18/09 & Température 18°C, point de rosée 9°C \\
Q1018 & QNH = 1018 hPa \\
\bottomrule
\end{tabular}

\begin{memobox}
Couverture~: FEW (1-2/8) $\to$ SCT (3-4/8) $\to$ BKN (5-7/8) $\to$ OVC (8/8).\\
\textbf{CAVOK} = Visi $>$10 km + pas de nuages $<$5000 ft + pas de phénomène significatif.
\end{memobox}

\newpage
%=============================================================
\section{AÉROPORTS \& INFRASTRUCTURE}
%=============================================================

\subsection{Paris-Charles de Gaulle (CDG / LFPG)}
\begin{itemize}[nosep]
\item \textbf{4 pistes}~: 08L/26R (4142m), 08R/26L (2700m), 09L/27R (2700m), 09R/27L (4200m)
\item Orientation~: Est-Ouest ($\sim$080°--090°)
\item Hub Air France depuis \textbf{1996}. $\sim$75 millions pax/an (2e en Europe)
\end{itemize}

\subsection{Paris-Orly (ORY / LFPO)}
\begin{itemize}[nosep]
\item \textbf{3 pistes}~: 02/20, 06/24, 07/25. Couvre-feu 23h30--06h00
\item $\sim$33 millions pax/an. Base principale Transavia France
\end{itemize}

\subsection{Codes OACI France}
\textbf{LF} = France. LFPG (CDG), LFPO (Orly), LFBO (Toulouse), LFLL (Lyon), LFML (Marseille), LFBD (Bordeaux), LFMN (Nice).

%=============================================================
\section{ALLIANCES \& FIDÉLITÉ}
%=============================================================

\subsection{SkyTeam --- 18 Membres (2025)}
Fondée le 22 juin 2000. Siège~: Amsterdam.

\begin{multicols}{3}
\begin{enumerate}[nosep]
\item \textbf{Air France}
\item \textbf{KLM}
\item \textbf{Delta Air Lines}
\item \textbf{Korean Air}
\item Aeroméxico
\item Aerolíneas Argentinas
\item Air Europa
\item China Airlines
\item China Eastern
\item Garuda Indonesia
\item Kenya Airways
\item Middle East Airlines
\item SAS
\item Saudia
\item TAROM
\item Vietnam Airlines
\item Virgin Atlantic
\item Xiamen Air
\end{enumerate}
\end{multicols}

\subsection{Les 3 Grandes Alliances}

\begin{tabular}{p{3cm}p{1.5cm}p{10cm}}
\toprule
\textbf{Alliance} & \textbf{Fondée} & \textbf{Compagnies phares} \\
\midrule
Star Alliance & 1997 & Lufthansa, United, ANA, Swiss, Turkish, Singapore Airlines \\
\textbf{SkyTeam} & \textbf{2000} & Air France, Delta, KLM, Korean Air \\
Oneworld & 1999 & British Airways, American, Qantas, JAL, Iberia, Cathay Pacific \\
\bottomrule
\end{tabular}

%=============================================================
\section{SÉCURITÉ \& ÉVÉNEMENTS MAJEURS}
%=============================================================

\begin{longtable}{p{2cm}p{3.5cm}p{9cm}}
\toprule
\textbf{Date} & \textbf{Événement} & \textbf{Détails} \\
\midrule
1994 & AF8969 détourné & GIA, intervention \textbf{GIGN} à Marseille-Marignane \\
\rowcolor{warnred}
\textbf{25/07/2000} & \textbf{Crash Concorde} & Gonesse (CDG). 113 morts. Débris pneu $\to$ réservoir percé $\to$ incendie \\
11/09/2001 & Attentats 11 sept. & 4 avions, $\sim$3000 morts. Révolution sûreté aérienne \\
\rowcolor{warnred}
\textbf{01/06/2009} & \textbf{Vol AF447} Rio-Paris & \textbf{A330}, 228 morts. \textbf{Sondes Pitot givrées} $\to$ perte contrôle $\to$ décrochage \\
15/01/2009 & Miracle Hudson & Cpt \textbf{Sullenberger}, A320, 155 survivants (tous) \\
2014 & MH370 disparu & B777, océan Indien, 239 à bord, jamais retrouvé \\
2015 & Germanwings 9525 & A320, crash volontaire Alpes, 150 morts \\
2018-19 & Crashes B737 MAX & Système \textbf{MCAS} défaillant, 346 morts, cloué au sol 20 mois \\
\bottomrule
\end{longtable}

%=============================================================
\section{CODES, COMMUNICATIONS \& RÉGLEMENTATION}
%=============================================================

\subsection{Codes Transpondeur d'Urgence}

\begin{tabular}{p{2cm}p{4cm}p{8.5cm}}
\toprule
\textbf{7500} & Détournement (hijack) & <<~75 = someone's taken alive~>> \\
\textbf{7600} & Panne radio & <<~76 = need to fix~>> \\
\textbf{7700} & Urgence générale & <<~77 = going to heaven~>> \\
\bottomrule
\end{tabular}

\subsection{Alphabet OACI}
A-Alpha, B-Bravo, C-Charlie, D-Delta, E-Echo, F-Foxtrot, G-Golf, H-Hotel, I-India, J-Juliet, K-Kilo, L-Lima, M-Mike, N-November, O-Oscar, P-Papa, Q-Quebec, R-Romeo, S-Sierra, T-Tango, U-Uniform, V-Victor, \textbf{W-Whiskey}, X-X-ray, Y-Yankee, Z-Zulu.

\subsection{Organismes Internationaux}

\begin{tabular}{p{2.5cm}p{6cm}p{6cm}}
\toprule
\textbf{OACI} (ICAO) & Org. Aviation Civile Internationale & Normes mondiales (1944, Convention de Chicago) \\
\textbf{IATA} & International Air Transport Association & Association des compagnies (1945) \\
\textbf{EASA} & EU Aviation Safety Agency & Sécurité aérienne européenne (2002) \\
\textbf{DGAC} & Direction Gén. Aviation Civile & Autorité française \\
\textbf{BEA} & Bureau Enquêtes et Analyses & Enquêtes accidents France \\
\textbf{Eurocontrol} & --- & Gestion trafic aérien européen \\
\bottomrule
\end{tabular}

\subsection{Licences de Pilote}

\begin{tabular}{p{2cm}p{5cm}p{7.5cm}}
\toprule
BIA & Brevet Initiation Aéro & Initiation théorique (lycée) \\
PPL & Private Pilot Licence & Vol privé, non rémunéré \\
CPL & Commercial Pilot Licence & Vol rémunéré, copilote \\
\rowcolor{aflight}
\textbf{ATPL} & Airline Transport Pilot Licence & \textbf{Commandant de bord} --- objectif des Cadets \\
MPL & Multi-crew Pilot Licence & Copilote en équipage multi-pilotes \\
\bottomrule
\end{tabular}

\subsection{Unités Aéronautiques}

\begin{tabular}{p{3cm}p{4cm}p{7.5cm}}
\toprule
Altitude & \textbf{Pieds (ft)} & 1 ft = 0.3048 m \\
Distance & \textbf{Mille nautique (NM)} & 1 NM = \textbf{1\,852 m} \\
Vitesse & \textbf{N\oe ud (kt)} & 1 kt = 1 NM/h = \textbf{1.852 km/h} \\
Pression & hPa & Standard~: 1013.25 hPa \\
\bottomrule
\end{tabular}

\newpage
%=============================================================
\section{QCM D'ENTRAÎNEMENT --- 100 QUESTIONS}
%=============================================================

\textit{Conditions réelles~: 15-30s par question. Bonne réponse~: +3 pts. Mauvaise~: $-$1 pt.}

% QCM PSY0 - 100 Questions
% Ce fichier est inclus par REVISION_PSY0.tex

\subsection*{Histoire (Q1--20)}

\begin{enumerate}[nosep]
\item Père de l'aviation en France ? \\ a) Wright \quad b) Blériot \quad c) Ader \quad d) Garros

\item Création d'Air France ? \\ a) 1929 \quad b) 1933 \quad c) 1936 \quad d) 1945

\item Première traversée de la Manche ? \\ a) Garros \quad b) Blériot \quad c) Lindbergh \quad d) Mermoz

\item Nombre de compagnies fusionnées pour créer AF ? \\ a) 3 \quad b) 4 \quad c) 5 \quad d) 6

\item Avion de Mermoz à sa disparition ? \\ a) Breguet XIV \quad b) Latécoère 300 \quad c) Spirit of St. Louis \quad d) Blériot XI

\item Premier avion de ligne à réaction ? \\ a) Caravelle \quad b) B707 \quad c) Comet \quad d) DC-8

\item Premier vol du Concorde ? \\ a) 1965 \quad b) 1969 \quad c) 1972 \quad d) 1976

\item Première traversée Atlantique Nord solo ? \\ a) Mermoz \quad b) Nungesser \quad c) Earhart \quad d) Lindbergh

\item Création d'Airbus ? \\ a) 1965 \quad b) 1968 \quad c) 1970 \quad d) 1975

\item Première femme brevetée pilote ? \\ a) Earhart \quad b) Bolland \quad c) Deroche \quad d) Boucher

\item Traversée de la Méditerranée 1913 ? \\ a) Blériot \quad b) Garros \quad c) Mermoz \quad d) Fabre

\item Convention de Chicago (1944) crée ? \\ a) IATA \quad b) OACI \quad c) EASA \quad d) Eurocontrol

\item Premier vol supersonique ? \\ a) Concorde \quad b) Yeager/X-1 \quad c) SR-71 \quad d) Armstrong

\item Hub CDG inauguré en ? \\ a) 1990 \quad b) 1993 \quad c) 1996 \quad d) 2000

\item Exploit d'Adrienne Bolland ? \\ a) Manche \quad b) Tour du monde \quad c) Andes \quad d) Méditerranée

\item Avion de Nungesser et Coli ? \\ a) Spirit of St. Louis \quad b) Éole \quad c) L'Oiseau Blanc \quad d) Croix du Sud

\item Premier vol en hydravion (1910) ? \\ a) Blériot \quad b) Garros \quad c) Fabre \quad d) Ader

\item Amelia Earhart solo Atlantique ? \\ a) 1928 \quad b) 1932 \quad c) 1935 \quad d) 1937

\item Mise en service Concorde chez AF ? \\ a) 1969 \quad b) 1972 \quad c) 1976 \quad d) 1980

\item SkyTeam fondée en ? \\ a) 1997 \quad b) 1999 \quad c) 2000 \quad d) 2003
\end{enumerate}

\subsection*{Air France \& Flotte (Q21--40)}

\begin{enumerate}[nosep, start=21]
\item DG groupe AF-KLM ? \\ a) Spinetta \quad b) Rigail \quad c) Smith \quad d) Janaillac

\item Emblème d'AF ? \\ a) Coq \quad b) Hippocampe ailé \quad c) Aigle \quad d) Fleur

\item Moteurs A330 ? \\ a) 1 \quad b) 2 \quad c) 3 \quad d) 4

\item Moteurs A340 ? \\ a) 2 \quad b) 3 \quad c) 4 \quad d) 6

\item Moteur du B777 ? \\ a) CFM56 \quad b) Trent \quad c) GE90 \quad d) PW4000

\item A350 en service chez AF depuis ? \\ a) 2015 \quad b) 2017 \quad c) 2019 \quad d) 2021

\item Retrait A380 chez AF ? \\ a) 2018 \quad b) 2019 \quad c) 2020 \quad d) 2022

\item Flotte principale Transavia ? \\ a) A320 \quad b) B737-800 \quad c) A220 \quad d) E190

\item A220 remplace ? \\ a) A321 \quad b) A318/A319 \quad c) B737 \quad d) A340

\item Slogan AF depuis 2022 ? \\ a) France is in the Air \quad b) S'envoler en toute élégance \quad c) Gagner le c\oe ur \quad d) Sky

\item Taille flotte AF ? \\ a) 150 \quad b) 200 \quad c) $\sim$228 \quad d) 350

\item A320 chez AF depuis ? \\ a) 1985 \quad b) 1988 \quad c) 1992 \quad d) 1995

\item Capacité A380 chez AF ? \\ a) 380 \quad b) 450 \quad c) 516 \quad d) 615

\item Vitesse croisière Concorde ? \\ a) M0.85 \quad b) M1.5 \quad c) M2.02 \quad d) M3

\item DG d'Air France (femme) ? \\ a) Couderc \quad b) Rigail \quad c) Bolland \quad d) Deroche

\item Anne Rigail nommée en ? \\ a) 2016 \quad b) 2017 \quad c) 2018 \quad d) 2020

\item Motoriste du CFM56 ? \\ a) RR \quad b) P\&W \quad c) CFM (GE+Safran) \quad d) GE seul

\item Vitesse croisière A350 ? \\ a) M0.78 \quad b) M0.82 \quad c) M0.85 \quad d) M0.92

\item Transavia créée en ? \\ a) 2000 \quad b) 2003 \quad c) 2006 \quad d) 2010

\item Nb de B777 chez AF (toutes versions) ? \\ a) $\sim$30 \quad b) $\sim$45 \quad c) $\sim$61 \quad d) $\sim$80
\end{enumerate}

\subsection*{Technique de Vol (Q41--60)}

\begin{enumerate}[nosep, start=41]
\item Le variomètre mesure ? \\ a) Vitesse \quad b) Altitude \quad c) Vitesse verticale \quad d) Cap

\item Le décrochage : toujours à la même ? \\ a) Vitesse \quad b) Altitude \quad c) Incidence \quad d) Masse

\item Gouverne du lacet ? \\ a) Ailerons \quad b) Profondeur \quad c) Direction \quad d) Volets

\item Assiette = ? \\ a) Portance+Traînée \quad b) Pente+Incidence \quad c) Alt+Vit \quad d) Cap+Route

\item Volets plus sortis en... \\ a) Décollage \quad b) Croisière \quad c) Atterrissage \quad d) Montée

\item Les spoilers ? \\ a) $\uparrow$portance \quad b) $\downarrow$portance + $\uparrow$traînée \quad c) Contrôle lacet \quad d) $\uparrow$vitesse

\item Code 7500 ? \\ a) Urgence \quad b) Radio \quad c) Détournement \quad d) Normal

\item Le TCAS est un système ? \\ a) Nav \quad b) Com \quad c) Anticollision \quad d) Météo

\item Ordre turboréacteur ? \\ a) Turb$\to$Comp$\to$Comb \quad b) Comp$\to$Comb$\to$Turb$\to$Tuyère \quad c) Comb$\to$Comp$\to$Turb \quad d) Tuyère$\to$Turb

\item Centrage avant = ? \\ a) Moins stable \quad b) Plus man\oe uvrant \quad c) Plus stable + plus conso \quad d) Plus rapide

\item V1 = vitesse de... ? \\ a) Rotation \quad b) Décision \quad c) Sécurité \quad d) Approche

\item Croisière A320 ? \\ a) M0.65 \quad b) M0.78 \quad c) M0.85 \quad d) M0.92

\item Reverses pour ? \\ a) Accélérer \quad b) Monter \quad c) Tourner \quad d) Freiner après atterrissage

\item Énergie gouvernes avion moderne ? \\ a) Électrique seul \quad b) Câbles \quad c) Hydraulique \quad d) Air comprimé

\item Prise statique bouchée $\to$ ? \\ a) Rien \quad b) GPS panne \quad c) Alti + vario bloqués \quad d) Moteurs

\item L'anémomètre utilise ? \\ a) Statique seule \quad b) Totale $-$ statique \quad c) Dynamique seule \quad d) GPS

\item Fly-by-wire inventé par ? \\ a) Boeing \quad b) Airbus (A320) \quad c) Dassault \quad d) Embraer

\item Virage 60° = facteur de charge ? \\ a) 1g \quad b) 1.4g \quad c) 2g \quad d) 3g

\item Les becs (slats) ? \\ a) $\uparrow$traînée seult \quad b) $\uparrow$incidence max, $\downarrow$Vs \quad c) $\downarrow$portance \quad d) Contrôle lacet

\item Taux de dilution CFM56 ? \\ a) $\sim$0.5 \quad b) $\sim$3 \quad c) $\sim$6 \quad d) $\sim$15
\end{enumerate}

\subsection*{Navigation \& Météo (Q61--80)}

\begin{enumerate}[nosep, start=61]
\item Pistes CDG ? \\ a) 2 \quad b) 3 \quad c) 4 \quad d) 5

\item Pistes Orly ? \\ a) 2 \quad b) 3 \quad c) 4 \quad d) 5

\item Température ISA au sol ? \\ a) 0°C \quad b) 10°C \quad c) 15°C \quad d) 20°C

\item Pression standard ? \\ a) 1000 \quad b) 1013.25 \quad c) 1025 \quad d) 1050 hPa

\item Gradient ISA ? \\ a) $-$1°C/1000ft \quad b) $-$2°C/1000ft \quad c) $-$3°C \quad d) $-$5°C

\item Temp ISA à 5\,000 ft ? \\ a) 10°C \quad b) 5°C \quad c) 0°C \quad d) $-$5°C

\item L'ILS fournit ? \\ a) Horizontal seul \quad b) Vertical seul \quad c) Les deux \quad d) Aucun

\item Membres SkyTeam 2025 ? \\ a) 15 \quad b) 17 \quad c) 18 \quad d) 20

\item Anticyclone (HN) tourne ? \\ a) Sens horaire \quad b) Antihoraire \quad c) Ne tourne pas \quad d) Variable

\item Jet stream se situe ? \\ a) Sol \quad b) 3000m \quad c) Tropopause \quad d) Stratosphère

\item Cb caractérisé par ? \\ a) Finesse \quad b) Bas \quad c) Grand dév. vertical \quad d) Pas de pluie

\item Ozone dans ? \\ a) Troposphère \quad b) Stratosphère \quad c) Mésosphère \quad d) Thermosphère

\item Vent le plus défavorable décollage ? \\ a) Face \quad b) Arrière \quad c) Travers \quad d) Calme

\item Tropopause à ? \\ a) 5 km \quad b) 8 km \quad c) 11 km \quad d) 15 km

\item Code OACI de CDG ? \\ a) LFPO \quad b) LFPG \quad c) LFBO \quad d) LFML

\item Djibouti est en ? \\ a) Afrique W \quad b) Afrique E \quad c) Océan Indien \quad d) Moyen-Orient

\item Luanda capitale de ? \\ a) Nigeria \quad b) Mozambique \quad c) Angola \quad d) Cameroun

\item Classe B en France ? \\ a) Contrôlée \quad b) Non contrôlée \quad c) Non utilisée \quad d) Militaire

\item 1 NM = ? \\ a) 1000 m \quad b) 1609 m \quad c) 1852 m \quad d) 2000 m

\item CAVOK signifie ? \\ a) Ciel couvert \quad b) Ceiling And Visibility OK \quad c) Clear Air \quad d) Pas de vent
\end{enumerate}

\subsection*{Divers \& Approfondissement (Q81--100)}

\begin{enumerate}[nosep, start=81]
\item W en alphabet OACI ? \\ a) Watt \quad b) William \quad c) Whiskey \quad d) Walter

\item Distance Paris-NY ? \\ a) 3500 \quad b) 4800 \quad c) 5850 \quad d) 7200 km

\item NY$\to$Paris plus rapide grâce au ? \\ a) Terre tourne \quad b) Jet stream \quad c) Avion léger \quad d) Route courte

\item Plus grande envergure ? \\ a) A380 \quad b) An-225 \quad c) Stratolaunch \quad d) B747

\item Flying Blue pour ? \\ a) AF seul \quad b) AF+KLM+Transavia+... \quad c) Tout SkyTeam \quad d) AF+Lufthansa

\item Code 7600 ? \\ a) Urgence \quad b) Détournement \quad c) Panne radio \quad d) Incendie

\item BEA enquête sur ? \\ a) Terrorisme \quad b) Accidents aériens \quad c) Météo \quad d) Passagers

\item ATPL = licence ? \\ a) Privé \quad b) ULM \quad c) Pilote de ligne \quad d) Planeur

\item Crash Gonesse 2000 ? \\ a) AF447 \quad b) Concorde \quad c) AF8969 \quad d) Germanwings

\item AF447 était un ? \\ a) B777 \quad b) A330 \quad c) A340 \quad d) B747

\item Miracle Hudson : Cpt ? \\ a) Piché \quad b) Sullenberger \quad c) Crespigny \quad d) Haynes

\item Anémomètre inventé par ? \\ a) Pitot \quad b) Badin \quad c) Blériot \quad d) Garros

\item Max 250 kt sous ? \\ a) FL50 \quad b) FL100 \quad c) FL200 \quad d) FL350

\item 1 kt = ? \\ a) 1 km/h \quad b) 1.5 km/h \quad c) 1.852 km/h \quad d) 2 km/h

\item Super Constellation ? \\ a) Biréacteur \quad b) Biplace \quad c) Quadrimoteur pistons \quad d) Monomoteur

\item Mustang = surnom du ? \\ a) F-16 \quad b) P-51 \quad c) Spitfire \quad d) Mirage

\item Haute pression au décollage = ? \\ a) Pire perf. \quad b) Meilleure perf. \quad c) Aucun effet \quad d) Plus de vent

\item QNH affiche ? \\ a) Hauteur \quad b) Altitude \quad c) FL \quad d) Température

\item Fondateurs SkyTeam ? \\ a) AF+BA+LH \quad b) AF+Delta+AeroMéxico+Korean \quad c) AF+KLM \quad d) AF+Delta+United

\item Mach 0.78 A320 = environ ? \\ a) 700 km/h \quad b) 828 km/h \quad c) 900 km/h \quad d) 1000 km/h
\end{enumerate}

\vspace{1cm}
\subsection*{Grille de Correction}

\begin{center}
\begin{tabular}{|c|c||c|c||c|c||c|c||c|c|}
\hline
\textbf{Q} & \textbf{R} & \textbf{Q} & \textbf{R} & \textbf{Q} & \textbf{R} & \textbf{Q} & \textbf{R} & \textbf{Q} & \textbf{R} \\
\hline
1 & c & 21 & c & 41 & c & 61 & c & 81 & c \\
2 & b & 22 & b & 42 & c & 62 & b & 82 & c \\
3 & b & 23 & b & 43 & c & 63 & c & 83 & b \\
4 & c & 24 & c & 44 & b & 64 & b & 84 & c \\
5 & b & 25 & c & 45 & c & 65 & b & 85 & b \\
6 & c & 26 & c & 46 & b & 66 & b & 86 & c \\
7 & b & 27 & c & 47 & c & 67 & c & 87 & b \\
8 & d & 28 & b & 48 & c & 68 & c & 88 & c \\
9 & c & 29 & b & 49 & b & 69 & a & 89 & b \\
10 & c & 30 & b & 50 & c & 70 & c & 90 & b \\
11 & b & 31 & c & 51 & b & 71 & c & 91 & b \\
12 & b & 32 & b & 52 & b & 72 & b & 92 & b \\
13 & b & 33 & c & 53 & d & 73 & b & 93 & b \\
14 & c & 34 & c & 54 & c & 74 & c & 94 & c \\
15 & c & 35 & b & 55 & c & 75 & b & 95 & c \\
16 & c & 36 & c & 56 & b & 76 & b & 96 & b \\
17 & c & 37 & c & 57 & b & 77 & c & 97 & b \\
18 & b & 38 & c & 58 & c & 78 & c & 98 & b \\
19 & c & 39 & c & 59 & b & 79 & c & 99 & b \\
20 & c & 40 & c & 60 & c & 80 & b & 100 & b \\
\hline
\end{tabular}
\end{center}


% Sections supplémentaires - inclus par REVISION_PSY0.tex
% ETOPS, CRM, Carburant, Stratégie AF, Systèmes avion, Balisage piste

%=============================================================
\section{ETOPS \& OPÉRATIONS ÉTENDUES}
%=============================================================

\subsection{Définition}
\textbf{ETOPS} = \textit{Extended-range Twin-engine Operational Performance Standards} (ou \textbf{EDTO} depuis 2012, OACI).\\
Règle qui permet aux \textbf{biréacteurs} d'opérer sur des routes éloignées d'un aéroport de déroutement. Le chiffre ETOPS = temps max \textbf{en minutes} avec \textbf{1 moteur en panne} pour rejoindre un aéroport.

\subsection{Niveaux de Certification}

\begin{tabular}{p{3cm}p{3cm}p{8.5cm}}
\toprule
\textbf{ETOPS} & \textbf{Couverture} & \textbf{Exemples} \\
\midrule
60 min & Routes courtes & Règle historique, avant ETOPS \\
120 min & Atlantique partiel & A300, premières certifications \\
\rowcolor{aflight}
\textbf{180 min (3h)} & \textbf{$\sim$95\% du globe} & A330, B767, B777, A320neo \\
240 min & 99\% du globe & B777-200LR \\
\rowcolor{aflight}
\textbf{370 min (6h10)} & \textbf{99.7\% du globe} & \textbf{A350-900} (1er certifié, 2014) \\
\bottomrule
\end{tabular}

\begin{memobox}
\textbf{Pourquoi l'ETOPS est important ?} Avant cette règle, les biréacteurs ne pouvaient pas traverser l'Atlantique~! Seuls les quadriréacteurs (A340, B747) le pouvaient. Grâce à l'ETOPS-180+, le B777 a rendu les quadriréacteurs obsolètes (plus chers, plus gourmands).
\end{memobox}

%=============================================================
\section{CARBURANT AVIATION}
%=============================================================

\subsection{Jet A-1 --- Fiche Technique}

\begin{tabular}{p{4.5cm}p{10cm}}
\toprule
\textbf{Type} & Kérosène aviation \\
Composition & Hydrocarbures C9--C16 (paraffines, aromatiques) \\
\textbf{Densité} & $\sim$\textbf{0.8 kg/L} (775--840 kg/m$^3$) \\
\textbf{Point d'éclair} & $>$ \textbf{38°C} (sécurité incendie) \\
\textbf{Point de congélation} & $\leq$ \textbf{$-$47°C} (vol en altitude) \\
Pouvoir calorifique & 43.15 MJ/kg \\
Ébullition & 150--250°C \\
Norme & DEF STAN 91-91, ASTM D1655 \\
\bottomrule
\end{tabular}

\subsection{SAF --- Carburant Aviation Durable}
\begin{itemize}[nosep]
\item SAF = \textit{Sustainable Aviation Fuel}, produit à partir de biomasse, huiles usagées, déchets
\item Réduit le CO$_2$ de 50--80\% sur le cycle de vie complet
\item Objectif AF~: \textbf{10\% de SAF} d'ici 2030
\item Mélangeable au Jet A-1 (max 50\% actuellement)
\item Pas de modification moteur nécessaire (``\textit{drop-in fuel}'')
\end{itemize}

\subsection{Avgas 100LL}
Carburant pour avions \textbf{à pistons} (avions légers). Essence contenant encore du plomb (coloration bleue). 100 octanes.

%=============================================================
\section{FACTEURS HUMAINS \& CRM}
%=============================================================

\subsection{CRM --- Crew Resource Management}

\begin{tabular}{p{4cm}p{10.5cm}}
\toprule
\textbf{Définition} & Formation aux compétences non-techniques pour améliorer la sécurité \\
Origine & Atelier NASA 1979, après crash Tenerife 1977 (583 morts) et UA 173 (1978) \\
Objectif & Gestion optimale de \textbf{toutes les ressources} disponibles (humaines, matérielles, informationnelles) \\
\bottomrule
\end{tabular}

\subsection{Compétences CRM}
\begin{enumerate}[nosep]
\item \textbf{Communication} --- phraséologie standard, read-back/hear-back
\item \textbf{Conscience situationnelle} --- modèle mental partagé de la situation
\item \textbf{Prise de décision} --- gestion du temps, évaluation des options
\item \textbf{Gestion de la charge de travail} --- priorisation, délégation
\item \textbf{Leadership et assertivité} --- le copilote \textbf{doit} dire s'il voit un problème
\item \textbf{Travail d'équipe} --- synergie PNT + PNC + ATC
\end{enumerate}

\subsection{Modèle de Reason --- ``Fromage Suisse''}

\begin{center}
\fbox{\parbox{13cm}{
Chaque \textbf{couche de défense} (procédures, formation, technologie, réglementation) est comme une tranche de fromage suisse avec des \textbf{trous} (faiblesses). Un accident survient quand les trous de \textbf{toutes les tranches s'alignent}, permettant au danger de traverser toutes les défenses.\\[0.3cm]
$\Rightarrow$ L'accident n'est \textbf{jamais} dû à une seule cause, mais à un \textbf{alignement} de facteurs (organisationnels + supervision + conditions + actes dangereux).
}}
\end{center}

\begin{memobox}
\textbf{80\% des accidents} en aviation sont liés aux \textbf{facteurs humains}. Le CRM a drastiquement réduit le taux d'accidents depuis les années 1980.
\end{memobox}

%=============================================================
\section{STRATÉGIE ENVIRONNEMENTALE --- AIR FRANCE ACT}
%=============================================================

\begin{tabular}{p{5cm}p{9.5cm}}
\toprule
\textbf{Programme} & \textbf{Air France ACT} \\
\midrule
Objectif 2030 (CO$_2$) & \textbf{$-$30\% par passager-km} vs 2019 (validé SBTi) \\
Objectif 2050 & Net zéro émissions \\
Incorporation SAF 2030 & \textbf{10\%} au niveau mondial \\
Flotte nouvelle génération & \textbf{70\%} d'ici 2030 (A220, A350, A320neo) \\
Investissement flotte & $>$1 Md\euro/an \\
Éco-pilotage & Montée continue, descente continue, roulage un moteur \\
Compensation carbone & Programme volontaire pour les passagers \\
\bottomrule
\end{tabular}

%=============================================================
\section{SYSTÈMES AVION}
%=============================================================

\subsection{Circuits Principaux d'un Avion de Ligne}

\begin{tabular}{p{3.5cm}p{4cm}p{7cm}}
\toprule
\textbf{Système} & \textbf{Source} & \textbf{Fonction} \\
\midrule
\textbf{Hydraulique} & 3 circuits (Vert/Bleu/Jaune sur Airbus) & Commandes de vol, train, freins, reverses \\
\textbf{Électrique} & Générateurs moteurs + APU + RAT (éolienne secours) & Instruments, éclairage, calculateurs, galley \\
\textbf{Pneumatique} & Prélèvement d'air moteurs (bleed air) & \textbf{Pressurisation cabine}, antigivrage, démarrage moteurs \\
\textbf{Carburant} & Réservoirs ailes + central & Alimentation moteurs, centrage (transfert) \\
\textbf{Oxygène} & Bouteilles (équipage) + générateurs chimiques (pax) & Urgence dépressurisation (masques tombent) \\
\bottomrule
\end{tabular}

\subsection{Pressurisation}
\begin{itemize}[nosep]
\item Cabine pressurisée à $\sim$\textbf{6\,000--8\,000 ft} d'altitude cabine (même si l'avion est à FL390)
\item Altitude cabine max certifiée~: 8\,000 ft (2\,438 m)
\item Au-dessus de \textbf{10\,000 ft} cabine (~3\,000 m)~: risque d'hypoxie $\Rightarrow$ masques à oxygène
\item \textbf{En cas de dépressurisation}~: descente d'urgence vers FL100 (10\,000 ft)
\end{itemize}

\subsection{APU (Auxiliary Power Unit)}
\begin{itemize}[nosep]
\item Petit turboréacteur auxiliaire dans la queue de l'avion
\item Fournit~: \textbf{électricité + air comprimé} au sol (sans moteurs principaux)
\item Permet le démarrage des moteurs principaux
\item Peut être utilisé en vol comme source de secours
\end{itemize}

%=============================================================
\section{BALISAGE ET MARQUAGE DES PISTES}
%=============================================================

\subsection{Numérotation des Pistes}
\begin{itemize}[nosep]
\item Le numéro = \textbf{cap magnétique arrondi à la dizaine} divisé par 10
\item Exemple~: piste orientée cap 086° $\Rightarrow$ piste \textbf{09}. L'autre sens~: 086°+180° = 266° $\Rightarrow$ piste \textbf{27}
\item Pistes parallèles~: suffixe L (Left), R (Right), C (Center)
\item CDG~: 08L/26R et 09R/27L (les 2 grandes pistes de 4\,200 m)
\end{itemize}

\subsection{Feux de Piste}

\begin{tabular}{p{4cm}p{10.5cm}}
\toprule
Feux de bord de piste & \textbf{Blancs} (jaunes sur les derniers 600m) \\
Feux de seuil & \textbf{Verts} (début de piste) \\
Feux de fin de piste & \textbf{Rouges} \\
Feux d'axe de piste & Blancs, puis alternés rouge/blanc, puis rouges \\
PAPI (indicateur de pente) & 4 feux~: 2 blancs + 2 rouges = sur le plan (3°) \\
\bottomrule
\end{tabular}

\begin{memobox}
\textbf{PAPI}~: 4 rouges = trop bas (danger !). 4 blancs = trop haut. 2 rouges + 2 blancs = correct (plan de 3°).
\end{memobox}

\subsection{Taxiways}
\begin{itemize}[nosep]
\item Ligne de roulage~: \textbf{jaune} (axe central)
\item Point d'arrêt avant piste~: \textbf{2 lignes jaunes continues + 2 tiretées} (NE JAMAIS franchir sans clairance ATC)
\end{itemize}

%=============================================================
\section{PHRASÉOLOGIE RADIO --- EXEMPLES}
%=============================================================

\subsection{Exemples de Communications}

\begin{tabular}{p{3cm}p{11.5cm}}
\toprule
\textbf{Phase} & \textbf{Exemple} \\
\midrule
Mise en route & \textit{``Air France 1234, request startup, information Bravo''} \\
Roulage & \textit{``Air France 1234, taxi to holding point runway 27R via Alpha, Bravo''} \\
Décollage & \textit{``Air France 1234, runway 27R, cleared for takeoff, wind 280/12''} \\
Montée & \textit{``Air France 1234, climb FL350''} \\
En route & \textit{``Air France 1234, maintain FL370, contact Paris Control 132.350''} \\
Descente & \textit{``Air France 1234, descend FL100, QNH 1018''} \\
Approche & \textit{``Air France 1234, cleared ILS approach runway 27R''} \\
Atterrissage & \textit{``Air France 1234, runway 27R, cleared to land, wind 260/08''} \\
\bottomrule
\end{tabular}

\subsection{Termes Importants}

\begin{tabular}{p{3cm}p{11.5cm}}
\toprule
\textbf{ROGER} & J'ai bien reçu votre message \\
\textbf{WILCO} & Will comply --- je vais exécuter \\
\textbf{AFFIRM} & Oui (ne pas dire ``yes'') \\
\textbf{NEGATIVE} & Non \\
\textbf{SAY AGAIN} & Répétez (ne pas dire ``repeat'', réservé à l'artillerie) \\
\textbf{MAYDAY} & Détresse (urgence vitale), répété 3 fois \\
\textbf{PAN PAN} & Urgence (non vitale), répété 3 fois \\
\textbf{SQUAWK} & Afficher un code transpondeur \\
\bottomrule
\end{tabular}

\begin{warnbox}
\textbf{MAYDAY} (3$\times$) = détresse vitale (feu, dépressurisation, panne moteur).\\
\textbf{PAN PAN} (3$\times$) = urgence non vitale (panne instrument, passager malade).\\
\textbf{Ne JAMAIS dire ``repeat''} en communication ATC~! On dit ``SAY AGAIN''.
\end{warnbox}

%=============================================================
\section{PERFORMANCES --- COMPLÉMENTS}
%=============================================================

\subsection{Effet du Vent sur les Performances}

\begin{tabular}{p{3cm}p{5.5cm}p{5.5cm}}
\toprule
\textbf{Type de vent} & \textbf{Au décollage} & \textbf{À l'atterrissage} \\
\midrule
\textbf{Vent de face} & $\downarrow$ distance de roulement (favorable) & $\downarrow$ distance d'atterrissage (favorable) \\
\rowcolor{warnred}
\textbf{Vent arrière} & $\uparrow$ distance de roulement (défavorable !) & $\uparrow$ distance d'atterrissage (défavorable !) \\
\textbf{Vent de travers} & Limitation max selon avion & Technique de décrabage \\
\bottomrule
\end{tabular}

\subsection{Effet de la Température et de l'Altitude}

\begin{tabular}{p{4cm}p{10.5cm}}
\toprule
\textbf{Température élevée} & Air moins dense $\Rightarrow$ moins de portance $\Rightarrow$ $\uparrow$ distance décollage, $\downarrow$ performances moteur \\
\textbf{Altitude terrain élevée} & Air moins dense (même effet). Aéroport en altitude = performances dégradées \\
\textbf{Altitude-densité} & Altitude fictive tenant compte de la vraie densité de l'air \\
\textbf{Piste mouillée} & $\uparrow$ distance de freinage. Risque d'aquaplaning \\
\bottomrule
\end{tabular}

\subsection{Distance Franchissable vs Autonomie}
\begin{itemize}[nosep]
\item \textbf{Distance franchissable}~: distance max en km (dépend de la vitesse, du vent, de la charge)
\item \textbf{Autonomie} (endurance)~: temps max en vol avec le carburant disponible
\item Le carburant embarqué inclut~: \textbf{trip fuel + réserves + déroutement + attente + taxi}
\item Réserve finale obligatoire~: \textbf{30 min} de vol (avion à réaction)
\end{itemize}

%=============================================================
\section*{QCM SUPPLÉMENTAIRES (Q101--120)}
\addcontentsline{toc}{section}{QCM SUPPLÉMENTAIRES (Q101--120)}
%=============================================================

\begin{enumerate}[nosep, start=101]
\item ETOPS-180 couvre quel \% du globe ? \\ a) 80\% \quad b) 90\% \quad c) 95\% \quad d) 100\%

\item Premier avion certifié ETOPS-370 ? \\ a) B777 \quad b) A350 \quad c) B787 \quad d) A330

\item Densité du Jet A-1 ? \\ a) 0.5 kg/L \quad b) 0.7 \quad c) $\sim$0.8 kg/L \quad d) 1.0

\item Point de congélation Jet A-1 ? \\ a) 0°C \quad b) $-$20°C \quad c) $-$38°C \quad d) $-$47°C

\item Le CRM a été développé après quel accident ? \\ a) Concorde \quad b) Tenerife 1977 \quad c) AF447 \quad d) 11 septembre

\item Modèle de Reason = ? \\ a) Dominos \quad b) Fromage suisse \quad c) Pyramide \quad d) Arbre

\item \% des accidents liés aux facteurs humains ? \\ a) 20\% \quad b) 50\% \quad c) $\sim$80\% \quad d) 100\%

\item Altitude cabine typique en croisière ? \\ a) 2\,000 ft \quad b) 6\,000--8\,000 ft \quad c) 15\,000 ft \quad d) 35\,000 ft

\item APU se trouve ? \\ a) Sous l'aile \quad b) Dans le cockpit \quad c) Dans la queue \quad d) Sous le fuselage

\item Numéro piste cap 270° ? \\ a) 09 \quad b) 27 \quad c) 18 \quad d) 36

\item PAPI 4 rouges = ? \\ a) Trop haut \quad b) Trop bas \quad c) Correct \quad d) Piste fermée

\item ``WILCO'' signifie ? \\ a) Roger \quad b) Will comply \quad c) Négatif \quad d) Répétez

\item MAYDAY = ? \\ a) Urgence \quad b) Détresse (vitale) \quad c) Panne radio \quad d) Information

\item SAF = ? \\ a) Safety Aviation Fuel \quad b) Sustainable Aviation Fuel \quad c) Standard Air Fuel \quad d) Secure Air Fuel

\item Objectif CO$_2$ AF 2030 ? \\ a) $-$10\% \quad b) $-$20\% \quad c) $-$30\% par pax-km \quad d) $-$50\%

\item Vent le plus favorable au décollage ? \\ a) Vent de face \quad b) Arrière \quad c) Travers \quad d) Aucun

\item Réserve finale min (jet) ? \\ a) 15 min \quad b) 30 min \quad c) 45 min \quad d) 60 min

\item Feux de seuil de piste = couleur ? \\ a) Rouge \quad b) Vert \quad c) Blanc \quad d) Jaune

\item Dépressurisation $\to$ descente vers ? \\ a) Sol \quad b) FL100 \quad c) FL200 \quad d) FL350

\item ``Say again'' remplace ? \\ a) Roger \quad b) Repeat \quad c) Affirm \quad d) Wilco
\end{enumerate}

\subsection*{Correction Q101--120}

\begin{center}
\begin{tabular}{|c|c||c|c||c|c||c|c|}
\hline
\textbf{Q} & \textbf{R} & \textbf{Q} & \textbf{R} & \textbf{Q} & \textbf{R} & \textbf{Q} & \textbf{R} \\
\hline
101 & c & 106 & b & 111 & b & 116 & a \\
102 & b & 107 & c & 112 & b & 117 & b \\
103 & c & 108 & b & 113 & b & 118 & b \\
104 & d & 109 & c & 114 & b & 119 & b \\
105 & b & 110 & b & 115 & c & 120 & b \\
\hline
\end{tabular}
\end{center}


\newpage
%=============================================================
\section*{FICHES MÉMO EXPRESS}
\addcontentsline{toc}{section}{FICHES MÉMO EXPRESS}
%=============================================================

\begin{tcolorbox}[colback=aflight, colframe=afblue, title=\textbf{Top 15 Dates Incontournables}]
\begin{multicols}{2}
\begin{enumerate}[nosep]
\item \textbf{1890} --- Ader (Éole)
\item \textbf{1903} --- Frères Wright
\item \textbf{1909} --- Blériot (Manche)
\item \textbf{1927} --- Lindbergh (Atlantique)
\item \textbf{1933} --- Création Air France
\item \textbf{1944} --- Convention Chicago $\to$ OACI
\item \textbf{1947} --- Mur du son (Yeager)
\item \textbf{1969} --- Concorde + B747 + Apollo 11
\item \textbf{1970} --- Création Airbus
\item \textbf{1988} --- A320 chez AF (fly-by-wire)
\item \textbf{1996} --- Hub CDG
\item \textbf{2000} --- SkyTeam
\item \textbf{2004} --- Air France-KLM
\item \textbf{2019} --- A350 chez AF
\item \textbf{2020} --- Retrait A380
\end{enumerate}
\end{multicols}
\end{tcolorbox}

\begin{tcolorbox}[colback=memoyellow, colframe=orange!70, title=\textbf{Chiffres à Connaître Par C\oe ur}]
\begin{multicols}{2}
\begin{itemize}[nosep]
\item Pistes CDG~: \textbf{4}
\item Pistes Orly~: \textbf{3}
\item Membres SkyTeam~: \textbf{18}
\item Flotte AF~: \textbf{$\sim$228 avions}
\item Croisière A320~: \textbf{Mach 0.78} (828 km/h)
\item Croisière A350~: \textbf{Mach 0.85}
\item Croisière Concorde~: \textbf{Mach 2.02}
\item ISA sol~: \textbf{15°C / 1013.25 hPa}
\item Gradient ISA~: \textbf{$-$2°C / 1000 ft}
\item Tropopause~: \textbf{11 km / 36000 ft}
\item 1 NM = \textbf{1\,852 m}
\item 1 kt = \textbf{1.852 km/h}
\item Max sous FL100~: \textbf{250 kt}
\item 1 hPa $\approx$ \textbf{28 ft}
\end{itemize}
\end{multicols}
\end{tcolorbox}

\begin{tcolorbox}[colback=okgreen, colframe=green!50!black, title=\textbf{Mnémotechniques}]
\begin{itemize}[nosep]
\item \textbf{Transpondeur}~: 75 (hijack) / 76 (radio) / 77 (emergency)
\item \textbf{Turboréacteur}~: Compresseur $\to$ Combustion $\to$ Turbine $\to$ Tuyère
\item \textbf{Axes}~: Roulis=Ailerons / Tangage=Profondeur / Lacet=Direction
\item \textbf{Calages}~: QNH=Altitude / QFE=Hauteur / 1013=FL
\item \textbf{Pression}~: <<~High to low, look out below!~>>
\item \textbf{Moteurs}~: GE90 $\to$ B777 | Trent $\to$ A350 | CFM56 $\to$ A320/A340/B737
\end{itemize}
\end{tcolorbox}

\vfill
\begin{center}
\rule{0.5\textwidth}{1pt}\\[0.5cm]
{\footnotesize Sources~: Annales PSY0 2018, 2019, 2020 (Pilotest/Aeronet), Fiches Culture G AF ASP v5.1, Air France Corporate 2025, DGAC, SkyTeam, Airbus, Boeing}\\[0.3cm]
{\large\textbf{Bonne chance pour ta sélection, futur(e) pilote~!}}
\end{center}

\end{document}
